\begin{abstract}

Die Digitalisierung der Terminverwaltung im veterinärmedizinischen Bereich weist im
Vergleich zur Humanmedizin noch erhebliche Lücken auf. Bestehende Softwarelösungen
decken entweder ausschließlich interne Praxisverwaltungsfunktionen ab oder bieten
einen zu breiten Funktionsumfang, der über die reine Terminverwaltung hinausgeht.
Im Rahmen dieser Arbeit wird daher untersucht, wie sich eine plattformübergreifende
Terminverwaltungsanwendung für Tierhalter im Kontext von Veterinärmedizinpraxen
konzipieren, umsetzen und hinsichtlich ihrer Funktionalität und Usability evaluieren
lässt.

Methodisch folgt die Arbeit dem Ansatz des Design Science Research nach Hevner et al.,
da im Mittelpunkt die Konstruktion eines Artefakts steht, das ein praxisrelevantes
Problem löst. Als Ergebnis wurde die mobile Anwendung \textit{petappoint} entwickelt.
Die Anwendung wurde mit React Native und Expo umgesetzt und ermöglicht Tierhaltern
die Registrierung, die Verwaltung von Haustierprofilen, die Suche nach
Tierarztpraxen sowie die Buchung und Stornierung von Terminen.

Zur Evaluation wurde eine aufgabenbasierte Nutzerstudie mit 20 Teilnehmenden
durchgeführt, die gleichmäßig auf iOS und Android aufgeteilt wurden. Alle
Teilnehmenden schlossen sämtliche Aufgaben erfolgreich ab, was einer Erfolgsrate
von 100\,\% entspricht. Der durchschnittliche System Usability Scale\,(SUS)-Score
von 94,34 entspricht der Note~A und belegt eine hohe wahrgenommene
Gebrauchstauglichkeit auf beiden Plattformen. Die Ergebnisse zeigen, dass eine
fokussierte, plattformübergreifende Terminverwaltungslösung im veterinärmedizinischen
Bereich sowohl technisch umsetzbar als auch von der Zielgruppe als intuitiv nutzbar
bewertet wird.

\end{abstract}

\cleardoublepage
\thispagestyle{empty}

\section*{Eigenständigkeitserklärung}

Ich, Aziz Erol, versichere hiermit, dass ich die vorliegende Arbeit selbstständig verfasst und keine anderen als die angegebenen Quellen und Hilfsmittel
benutzt habe, wobei ich alle wörtlichen und sinngemäßen Zitate als solche gekennzeichnet habe.

Es wurden folgende KI-Tools verwendet: Claude Sonnet 4.6 (Anthropic), Perplexity Sonar 2 (Perplexity AI) sowie DeepL (DeepL SE). Etwaige durch die KI vorgenommenen Änderungen
und Empfehlungen wurden von mir kritisch geprüft und gegebenenfalls angepasst, um die wissenschaftliche Integrität und die Originalität der Arbeit sicherzustellen.

Diese Arbeit wurde bisher in gleicher oder ähnlicher Form keiner anderen  Prüfungsbehörde vorgelegt und auch nicht veröffentlicht.

\bigskip
Berlin, den \today

\bigskip
\bigskip
\noindent\rule{6cm}{0.4pt}

\addtocounter{page}{-1}
\cleardoublepage











